% A Minimal Epistemic Core -- LuaLaTeX.  Build: sh build.sh  (or: lualatex epistemic_core.tex, twice)
\documentclass[11pt,a4paper,oneside]{scrartcl}

% ---- page and type ------------------------------------------------------------
\usepackage[textwidth=29pc,textheight=46\baselineskip,top=1.15in,headsep=2.2\baselineskip,
            footskip=3\baselineskip,heightrounded]{geometry}
\usepackage{amsmath}
\usepackage{unicode-math}
\setmainfont{Libertinus Serif}[Numbers=OldStyle]
\setsansfont{Libertinus Sans}
\setmathfont{Libertinus Math}
\usepackage[babel=true]{microtype}
\usepackage[english]{babel}
\usepackage[autostyle]{csquotes}
\usepackage{booktabs,tabularx,array}
\usepackage{flafter} % floats never precede their place in the text
\usepackage{tikz}
\usetikzlibrary{arrows.meta}
% Figure conventions: terms are marks (positions), never areas; F and R are column headings (roles),
% never enclosures; one arrow style is reserved for receptive dependence; grey marks what belongs
% only to the theorist's description.
\colorlet{desc}{black!45}
\tikzset{
  term/.style  = {circle, fill, inner sep=0pt, minimum size=3.4pt},
  rel/.style   = {line width=.5pt},
  corr/.style  = {densely dotted, line width=.6pt, desc},
  dep/.style   = {-{Stealth[length=2.6mm,width=2mm]}, line width=.8pt},
  side/.style  = {font=\scriptsize\scshape, text=desc},
  note/.style  = {font=\footnotesize, text=desc, align=center},
  panel/.style = {font=\footnotesize},
}
\usepackage[backend=biber,style=authoryear,maxcitenames=2,giveninits=true,uniquename=false,
            doi=false,url=false,isbn=false,dashed=false]{biblatex}
\addbibresource{refs.bib}
\renewcommand*{\bibfont}{\small}
\setlength{\bibhang}{1.4em}
\renewbibmacro{in:}{\ifentrytype{article}{}{\printtext{\bibstring{in}\intitlepunct}}}
\usepackage{scrlayer-scrpage}
\deffootnote[1.4em]{1.4em}{1em}{\makebox[1.4em][l]{\thefootnotemark.}}
\setlength{\footnotesep}{.7\baselineskip}
\setkomafont{footnote}{\small}
\setkomafont{caption}{\small}
\setkomafont{captionlabel}{\scshape}
\setcapindent{0pt}
\renewcommand*{\captionformat}{.\enskip}
\usepackage[hidelinks,pdfusetitle]{hyperref}

\linespread{1.06}
\setlength{\parindent}{1.4em}
\widowpenalty=4000 \clubpenalty=4000 \interfootnotelinepenalty=0
\setlength{\emergencystretch}{1.5em}

% ---- running heads --------------------------------------------------------------
\newcommand{\shorttitle}{A Minimal Epistemic Core}
\newcommand{\theauthor}{Mario Asis}
\clearpairofpagestyles
\ohead{\textsc{\MakeLowercase{\shorttitle}}}
\ihead{\textsc{\MakeLowercase{\theauthor}}}
\cfoot*{\pagemark}
\setkomafont{pageheadfoot}{\normalfont\footnotesize}
\setkomafont{pagenumber}{\normalfont\small}
\pagestyle{scrheadings}

% ---- headings ---------------------------------------------------------------------
\setkomafont{disposition}{\normalfont}
\setkomafont{section}{\normalfont\large\scshape}
\RedeclareSectionCommand[beforeskip=-2.2\baselineskip plus -.4\baselineskip, % negative: no indent after heading
                         afterskip=.9\baselineskip]{section}
\renewcommand*{\sectionformat}{\thesection.\enskip}

% ---- displayed commitments and definitions -------------------------------------
% Indented block with a small-capitals label; the statement itself is set in italics.
\makeatletter
% Optional argument: the tag by which the text refers back to the statement, e.g. (D).
\newenvironment{thesis}[2][]
  {\@beginparpenalty=\@M % keep the announcing sentence with the statement
   \begin{list}{}{\setlength{\leftmargin}{2.2em}\setlength{\rightmargin}{2.2em}
                  \setlength{\topsep}{.9\baselineskip plus .2\baselineskip}}%
   \item\relax{\upshape\if\relax\detokenize{#1}\relax\else(#1)\enspace\fi\scshape\MakeLowercase{#2}}\par\nopagebreak\vspace{.2\baselineskip}\itshape\ignorespaces}
  {\end{list}}
\makeatother

% ---- notation ----------------------------------------------------------------------
\newcommand{\F}{\ensuremath{F}}
\newcommand{\R}{\ensuremath{R}}
\newcommand{\C}{\ensuremath{C}}

\hypersetup{pdftitle={A Minimal Epistemic Core},pdfauthor={Mario Asis},
            pdfsubject={Necessary conditions of any epistemic relation: differentiation, receptive dependence, and registration}}

% ===================================================================================
\begin{document}

\thispagestyle{plain.scrheadings}
\begin{center}
  \vspace*{.2\baselineskip}
  {\LARGE A Minimal Epistemic Core\par}
  \vspace{.7\baselineskip}
  {\large\itshape Differentiation, Receptive Dependence, and Registration\par}
  \vspace{1.6\baselineskip}
  {\scshape \MakeLowercase{Mario Asis}\par}
  \vspace{1.8\baselineskip}
\end{center}

\begin{center}
\begin{minipage}{24pc}
  \small
  \noindent{\scshape abstract.}\enspace
  This paper proposes a minimal core of necessary conditions for any epistemic relation: conditions
  without which nothing could be perceived, registered, or claimed about anything, whatever the
  domain in which the relation is realized. The core consists of three commitments, none of which can
  be derived from the others: that the field under consideration admits differentiation; that there
  is a second relatum admitting differentiation of its own; and that a relation of receptive
  dependence matches a contrast realized in the field with a contrast in the second relatum, the
  latter obtaining partly in virtue of the former. Contrasts are throughout between actual terms,
  never between alternatives, so that the core requires no repertoire of possibilities on either
  side. On this basis, registration is defined in a strictly structural sense that presupposes
  neither understanding, nor representation, nor order. The commitments are formal: they presuppose
  no spatial, temporal, or physical structure, and they are claimed to be necessary, not sufficient.
  The paper further shows that target and receiver are positions in the relation of dependence
  rather than natures of its relata, that specifying an articulation selects it without creating
  it, and that acquisition and retention require an ordering which the core does not contain and
  whose kind it leaves open.
\end{minipage}
\end{center}
\vspace{1.2\baselineskip}

% -----------------------------------------------------------------------------------
\section{Introduction}\label{sec:intro}

Theories of knowledge ordinarily begin with a subject who believes, judges, or is justified, and ask
under what conditions such states amount to knowledge. This paper addresses a prior question: what
must be admitted for any epistemic relation to be possible at all---for anything to be perceived,
registered, or claimed about anything else. Its answer is a set of necessary conditions, not a
sufficient account. It identifies what any theory of perception, belief, or knowledge presupposes,
and leaves to such theories what makes a relation epistemic in the full sense. The assessment of
judgement, in particular, lies outside its scope.

The guiding premise is that an epistemic relation is always a relation to something in particular.
To perceive or claim anything about $X$, $X$ must be distinguishable from what it is not, and that
distinguishing must depend on $X$; otherwise what is perceived or claimed would not be about $X$.
The commitments introduced below are what this premise requires when no further assumption is made:
a differentiation in what is concerned, a differentiation elsewhere, and a dependence of the second
on the first.

The conditions are formal. They are stated without reference to space, time, matter, or causation,
and they are meant to hold whatever the domain in which an epistemic relation is realized, whether
perceptual, physical, linguistic, or mathematical. Examples drawn from particular domains show how
the conditions are met there; they do not fix what the conditions are. For the same reason,
necessity is claimed only where it can be argued for, and features of the presentation that are
conventions rather than commitments are marked as such.

The reconstruction proceeds by minimal commitment. Each condition is introduced only at the point
where the account cannot proceed without it, and each is accompanied by a statement of what it does
not establish. One methodological caution applies throughout. The exposition necessarily employs
ordinary logical resources---identity, distinction, relation, and inference---but using them to state
the commitments does not attribute them to what the commitments describe, and in particular not to a
knowing subject. A distinction drawn in the exposition is not thereby a distinction drawn by its
object.

The proposal has close neighbours. \textcite{dretske1981} analyses the information carried by a
signal in terms of nomic dependence and conditional probability, and separates it from meaning and
belief. Floridi's diaphoric definition of data takes a lack of uniformity---a difference---as the
most elementary constituent of information \autocite{floridi2011}, and his informational structural
realism holds that the relation of difference is constitutive of its relata \autocite{floridi2008}.
Smith's metaphysics of presence makes \emph{registration} its central notion, as the way in which a
world comes to be taken as one of objects \autocite{smith1996}.\footnote{The present use of the term
is narrower. Smith's registration is performed by subjects and yields objects; registration here
presupposes neither and concerns only the dependence of one contrast on another.} Enactivist accounts describe
cognition as the structural coupling of a system with a medium, both of which belong to a larger
organization \autocite{maturana1987,thompson2007}. The present account differs from each in three
respects. It separates the differentiability of what is known, the differentiability of what
receives it, and the dependence that connects them, and treats each as an independent commitment. It
states that dependence between contrasts of actual terms, without probabilities, causes, or
alternatives on either side, so that it applies beyond the physical domain and requires no
repertoire of possibilities. And it distinguishes the direction of dependence from any order of
occurrences, leaving open what kind of order further epistemic operations require.

Section~\ref{sec:diff} introduces the differentiability of the field, and Section~\ref{sec:art}
distinguishes that capacity from any particular articulation. Section~\ref{sec:recv} introduces a
second relatum, and Section~\ref{sec:dep} the relation of receptive dependence that connects the two
and defines their roles. Section~\ref{sec:reg} defines registration on that basis and states what it
does not include, and Section~\ref{sec:order} shows that the resulting structure contains no order.
Section~\ref{sec:concl} summarizes the commitments, the limits of what they establish, what they
exclude, and the questions they leave open.

% -----------------------------------------------------------------------------------
\section{Differentiability}\label{sec:diff}

Let \F{} designate the \emph{field}: whatever is under consideration. The term fixes the scope of the
inquiry; it implies neither a spatial region, nor a collection of discrete objects, nor something
already presented to a subject. Nor does the notation establish a boundary intrinsic to reality:
\F{} identifies what the account concerns, and any claim about its actual boundaries would require
separate justification. The first commitment is the following.

\begin{thesis}[D]{Differentiability}
  \F{} admits differentiation, although no particular differentiation is thereby realized or
  specified.
\end{thesis}

Thesis (D) is stronger than the claim that different descriptions of \F{} can be devised. It
attributes to \F{} a capacity to enter into a differentiated articulation, so that differentiation is
not wholly attributable to the theorist's notation. At the same time, it is weaker than any claim
about which articulation is possible: it leaves open whether some articulation is uniquely
privileged, and whether its terms would be objects, qualities, regions, values, or something else.

(D) is accordingly a modal commitment; it concerns what \F{} admits. The capacity it attributes
belongs to \F{} itself. It is not a stock of possible configurations held anywhere, whether in a mind
or in some other repository, and it implies nothing about change or about any order, since the
possibility of differentiating something is distinct from the possibility of its changing.

(D) is not deduced from the mere existence of \F, for the fact that something is under consideration
establishes nothing about the structure it possesses. It is admitted on the following ground. A
rival view holds that all differentiation is drawn by whatever receives it: that distinctions are
made, not found.\footnote{The view has many forms. Spencer-Brown's calculus of indications begins
from the injunction to draw a distinction \autocite{spencerbrown1969}, which Luhmann makes the basis
of a constructivist theory of observation \autocite{luhmann1995}; Kant holds that the
combination of a manifold can never come to us through the senses \autocite[B129--130]{kant1998},
although for him a manifold is still given; and enactivist accounts treat the significance of a
difference as brought forth by the system for which it is significant \autocite{thompson2007}.} The
present account does not deny that a receiver contributes differentiations of its own;
Section~\ref{sec:recv} requires that it do so. But if every differentiation were supplied by the
receiver, nothing in \F{} would correspond to it, and the registration of a difference in \F{} could
not be distinguished from the receiver's merely differing within itself. By the guiding premise, what
was registered would then not be about \F. (D) is the condition under which that distinction can be
drawn. It is therefore a necessary condition of the explanation pursued here, not a sufficient
account of epistemic possibility. It also involves no knowing subject: that \F{} admits
differentiation does not imply that anything differentiates it.

% -----------------------------------------------------------------------------------
\section{Differentiability and Articulation}\label{sec:art}

Differentiability must be distinguished from the realization of any particular differentiation:
that \F{} admits differentiation does not entail that a differentiation has been
realized.\footnote{The refusal to ground a difference in terms identified before it has a precedent
in Deleuze's critique of the subordination of difference to identity \autocite[ch.~1]{deleuze1994},
and his distinction between \emph{differentiation} and \emph{differenciation}---the determination of
a structure and its actualization (ch.~4)---is akin to the distinction drawn in this section.
Deleuze also rejects the possible as a mere image of the real; the present account shares the
refusal of a stock of possibilities, although it retains, in (D), a modal claim about what \F{}
admits.}

A particular differentiation of \F{} will be called an \emph{articulation} and written
\[
  \langle a,\delta,b\rangle_F ,
\]
where $a$ and $b$ are its \emph{terms} and $\delta$ is its \emph{contrast}, the respect in which the
terms differ. An articulation is \emph{realized} when \F{} exhibits it, that is, when its terms are
actual and differ in the respect $\delta$. The terms of a realized articulation are both actual. They
are not alternatives of which one obtains rather than the other, and the contrast is not a choice
between them; it is the difference between two terms that are both there. The schema has two terms
for simplicity; the account requires only that an articulation have at least two. Nor does the
schema presuppose that $a$ and $b$ exist prior to the articulation as independent objects; it marks
positions within the articulation described.

Two misreadings of the schema should be excluded. First, $\delta$ is not a third item situated
between $a$ and $b$. To treat it as a separator would presuppose spatial organization, whereas
$\delta$ expresses the differentiation itself. Second, $b$ need not comprise everything that remains
once $a$ has been selected. That reading would require an exhaustive partition of \F, and a contrast
between two terms is not equivalent to a division into \enquote{this} and \enquote{everything else
within the field}: the latter adds the requirement that every case within the scope fall on one
side or the other. Figure~\ref{fig:contrast} sets the two constructions side by side.

\begin{figure}[tp]
  \centering\small
  \begin{tikzpicture}
    % (a) a contrast: two positions and the respect in which they differ
    \node[term, label=below:$a$] (a) at (0,0) {};
    \node[term, label=below:$b$] (b) at (2.2,0) {};
    \draw[rel] (a) to[bend left=55] node[above] {$\delta$} (b);
    \node[panel] at (1.1,-1.35) {(a) contrast};
    % (b) an exhaustive partition, deliberately drawn as a region
    \begin{scope}[xshift=4.1cm]
      \draw[desc, fill=desc!12, rounded corners=5pt] (0,-0.8) rectangle (4.4,0.9);
      \draw[desc] (1.4,-0.8) -- (1.4,0.9);
      \node[term, label=below:$a$] at (0.7,-0.05) {};
      \node[note] at (0.7,0.5) {this};
      \node[note] at (2.9,0.05) {everything else\\[-1pt]within the field};
      \node[panel] at (2.2,-1.35) {(b) exhaustive partition};
    \end{scope}
  \end{tikzpicture}
  \caption{Two constructions over \F. (a)~An articulation $\langle a,\delta,b\rangle_F$: $a$ and $b$
    mark positions, both actual, and the contrast $\delta$, drawn as a relation rather than as an
    item between them, is the respect in which they differ. (b)~An exhaustive partition, which adds
    the requirement that every case fall on one side or the other. Only (b) is drawn as a region,
    because only (b) carries the extensional commitment that a region suggests; in (a) the
    placement of the marks has no spatial meaning.}
  \label{fig:contrast}
\end{figure}

An articulation is \emph{specified} when its terms and contrast are fixed, and it is the account, not
\F, that specifies it. This might seem to reintroduce the dependence on notation that (D) excludes.
It does not. Specifying an articulation selects one of the differentiations that \F{} is claimed to
admit; it does not make \F{} admit it, still less realize it. Whether \F{} realizes the specified
articulation, and whether anything depends on it, are not settled by the specification, which can
therefore fail. Specification fixes what an inquiry concerns, as \F{} fixes its scope; it does not
supply what the inquiry finds.

These qualifications mark the first limit of the reconstruction. Differentiability permits
articulation, but it neither produces an actual articulation nor determines its form. An account of
any particular case must therefore specify the articulation it concerns; it cannot derive that
articulation from possibility alone. Moreover, a realized articulation remains an articulation of
\F, and nothing established so far implies that it is registered anywhere else.

% -----------------------------------------------------------------------------------
\section{The Second Relatum}\label{sec:recv}

For an articulation of \F{} to be registered, a differentiation must obtain somewhere other than in
the articulation itself. Let \R{} designate a second relatum, distinct from \F, in which it
obtains.\footnote{The vocabulary of reception recalls Kant's characterization of sensibility as
receptivity, the capacity to acquire representations through being affected by objects
\autocite[A19/B33; cf.\ A50/B74]{kant1998}. The receiver introduced here is thinner: it involves no
representations, and nothing is claimed about forms of intuition.} \R{} is not a person, a conscious
subject, or an independently bounded system, and it may ultimately prove to be an aspect of a larger
organization that also includes \F, since a distinction of relata does not entail a distinction of
substances.

\R{} must admit differentiation of its own. Otherwise the account could describe a contrast in \F{}
but nothing in \R{} that would distinguish its registration from its non-registration. The
requirement concerns the organization of \R, not the possession of concepts or classifications. The
second commitment is therefore the following.

\begin{thesis}[R]{Second relatum}
  There is a second relatum \R, distinct from \F, that admits differentiation of its own.
\end{thesis}

The terms of a realized differentiation of \R{} will be called its \emph{conditions} and written
$r, r', \dots$. As with articulations of \F, the conditions of a realized differentiation are both
actual: $r$ and $r'$ are not alternatives between which \R{} is poised, but distinct conditions that
\R{} exhibits. Nothing requires that a differentiation of \R{} have the same form as an articulation
of \F; the notation names its terms without assigning them any organization.

That (R) is a separate commitment can be shown without a theory of cognition. A contrast in \F{} and
a contrast between conditions of \R{} are not the same fact: an articulation of \F{} may be realized
while nothing in \R{} answers to it, and \R{} may exhibit contrasts that answer to nothing in \F. (R)
therefore does not follow from (D); it becomes necessary only when the account passes from
differentiation to registration.

(R) also fixes how much individuation the account admits. There must be enough to distinguish \F{}
from \R, and to take $r$ and $r'$ as conditions of one and the same relatum. There need not be
enough to identify enduring objects or successive states of a continuing subject: $r$ and $r'$ are
co-present conditions of \R, neither alternatives nor stages in its history. To treat these stronger
forms of individuation as already established would build into the foundation a persisting, bounded
bearer---which is precisely what an account of a knowing subject must eventually explain.

% -----------------------------------------------------------------------------------
\section{Receptive Dependence}\label{sec:dep}

Theses (D) and (R) leave \F{} and \R{} unconnected, for a field that admits differentiation and a
second relatum that admits differentiation of its own may be entirely unrelated. A further,
relational commitment is therefore required. Let \C{} designate \emph{receptive dependence}.

\begin{thesis}[C]{Receptive dependence}
  Through \C, the contrast of a specified articulation realized in \F{} is matched by a contrast
  between conditions of \R, and the latter contrast obtains partly in virtue of the former.
\end{thesis}

The thesis has two clauses, and both are needed. The first, \emph{matching}, requires that the terms
of the articulation be paired with conditions of \R{} that differ: $a$ with $r$ and $b$ with $r'$,
where $r$ and $r'$ are distinct. Were every term paired with the same condition, nothing in \R{}
would answer to the contrast $\delta$. Matching alone, however, is not dependence. Two structures may
resemble one another without either registering the other,\footnote{\textcite[ch.~1]{goodman1976}
makes the parallel point about representation: resemblance is reflexive and symmetric, whereas
representation is neither, so resemblance cannot be what representation consists in.} and two
contrasts may correspond term by term without either obtaining because of the other
(Figure~\ref{fig:correspondence}). Matching is symmetric: a correspondence that pairs $a$ with $r$
pairs $r$ with $a$, and can be read from \R{} to \F{} as well as from \F{} to \R.

\begin{figure}[tp]
  \centering\small
  \begin{tikzpicture}
    \foreach \x/\lab/\links in {0/{(a) resemblance}/0, 5.4/{(b) correspondence}/1} {
      \begin{scope}[xshift=\x cm]
        \node[side] at (0,1.25) {field (\F)};
        \node[side] at (2.4,1.25) {second relatum (\R)};
        \node[term, label=above:$a$]  (a) at (0,0.5)  {};
        \node[term, label=below:$b$]  (b) at (0,-0.5) {};
        \node[term, label=above:$r$]  (r) at (2.4,0.5)  {};
        \node[term, label=below:$r'$] (s) at (2.4,-0.5) {};
        \draw[rel] (a) to[bend right=60] node[left] {$\delta$} (b);
        \draw[rel] (r) to[bend left=60] (s);
        \ifnum\links=1
          \draw[corr] (a) -- (r);
          \draw[corr] (b) -- (s);
        \fi
        \node[panel] at (1.2,-1.55) {\lab};
      \end{scope}
    }
  \end{tikzpicture}
  \caption{Two relations that fall short of receptive dependence. (a)~Resemblance: a contrast in
    \F{} and a contrast between conditions of \R{} share a form, but neither registers the other.
    (b)~Correspondence: $a$ is paired with $r$, and $b$ with $r'$. The pairing satisfies the first
    clause of~(C) but is drawn without direction, since it can be read either way. Neither panel
    establishes~(C).}
  \label{fig:correspondence}
\end{figure}

The second clause supplies what matching lacks. The contrast in \R{} obtains \emph{in virtue of} the
contrast in \F: because of it, and not conversely or merely alongside it. The phrase is not meant to
identify \C{} with metaphysical ground; it borrows only two features that the literature on ground
has made familiar, namely that dependence of this kind is asymmetric and that it relates what
actually obtains, without appeal to what would obtain otherwise \autocite{fine2012}. (C) is
accordingly not a counterfactual claim. It does not say what \R{} would be if \F{} were otherwise; it
relates two contrasts that are both realized.\footnote{Counterfactual analyses of dependence require
alternatives to the actual case. Where contrasts are necessary, as in mathematics, they meet a
further difficulty: on the standard semantics, counterfactuals with impossible antecedents are all
vacuously true \autocite{lewis1973}; for an alternative semantics, see \textcite{nolan1997}. Neither
difficulty affects (C).}

The dependence is partial. Which condition of \R{} answers to which term depends also on the
manner in which \R{} is receptive, that is, on the organization of \R. The clause requires only that
the contrast in \R{} obtain \emph{partly} in virtue of the contrast in \F, the remainder being
supplied by \R{} itself.\footnote{This accommodates the enactivist point that a perturbation
triggers, but does not specify, the changes it occasions in a system, which are fixed by the
system's own structure \autocite{maturana1987}.} The point is captured in Bateson's definition of
information as \enquote{a difference which makes a difference} \autocite{bateson1972}, anticipated in
MacKay's \enquote{a distinction which makes a difference} \autocite{mackay1969}, which
compresses into one phrase what the present account keeps apart: the first difference belongs to
\F, the second to \R, and the making is the dependence \C. For the same reason, the content of (C)
should not be concealed behind appeals to \enquote{systematic} or \enquote{relevant} coupling, which
import what the argument has not supplied, namely a standard of relevance or a regularity across
cases.

The asymmetry of \C{} defines the roles of its relata. In an instance of \C, the \emph{target} is the
relatum whose contrast is depended upon, and the \emph{receiver} is the relatum whose contrast
depends. Roles are thus positions in a relation, not natures of relata. The same relatum may be
target in one instance of \C{} and receiver in another, and two relata may each depend on the other
in different respects. This also shows that the account is not circular. Neither (R) nor the
definition of \C{} uses the notion of a receiver: \C{} is introduced as an asymmetric relation between
contrasts, and the roles are read off it. \F{} and \R{} were named in anticipation of the positions
they occupy in the instances considered here.

Two limitations of (C) should be noted. First, the dependence is admitted, not derived from the
existence of differentiated relata; a relation cannot be obtained merely by describing its relata.
\C{} is introduced because registration requires a connection that neither relatum supplies on its
own. Second, the asymmetry of \C{} is primitive in the core. What grounds it in a particular domain
is a question for that domain: in physical domains it may be causal, in others it may be logical or
constitutive. The core requires that there be such an asymmetry, not that it be of any one kind.

% -----------------------------------------------------------------------------------
\section{Registration}\label{sec:reg}

The three commitments permit a definition of registration in a restricted, structural sense.

\begin{thesis}{Definition of registration}
  A specified articulation realized in \F{} is registered by \R{} when \C{} holds between them: its
  contrast is matched by a contrast between conditions of \R, which obtains partly in virtue of it.
\end{thesis}

\begin{figure}[tp]
  \centering\small
  \begin{tikzpicture}
    \node[side] at (0,1.45) {target (\F)};
    \node[side] at (5.0,1.45) {receiver (\R)};
    % target-side articulation
    \node[term, label=above:$a$] (a) at (0,0.6)  {};
    \node[term, label=below:$b$] (b) at (0,-0.6) {};
    \draw[rel] (a) to[bend left=55] coordinate[pos=.5] (dm) (b);
    \node[anchor=east, inner sep=2pt] at (dm) {$\delta$};
    % receiving contrast, with a different organization
    \node[term, label=below:$r$]  (r) at (4.3,-0.45) {};
    \node[term, label=below:$r'$] (s) at (5.7,-0.45) {};
    \draw[rel] (r) to[bend left=55] coordinate[pos=.5] (rm) (s);
    % receptive dependence
    \draw[dep] ([xshift=3pt]dm) -- node[above, pos=.5] {$C$} ([yshift=5pt]rm);
    % the three contributions
    \node[note] at (0,-1.65)   {realized\\[-1pt]articulation};
    \node[note] at (2.5,-1.65) {receptive\\[-1pt]dependence};
    \node[note] at (5.0,-1.65) {receiving\\[-1pt]contrast};
    \draw[desc, line width=.4pt] (-1.4,-2.25) -- (6.4,-2.25);
    \node[note, below] at (2.5,-2.25) {relative to the manner in which \R{} is receptive};
  \end{tikzpicture}
  \caption{Registration. Through \C, the contrast $\delta$ of the realized articulation
    $\langle a,\delta,b\rangle_F$ is matched by the contrast between the conditions $r$ and $r'$ of
    \R, which obtains partly in virtue of it. All four terms are actual. The receiving contrast is
    drawn with a different organization, since it need neither resemble nor reproduce the
    articulation. The arrow denotes receptive dependence only: it represents neither a transfer nor
    an order of occurrences, and the columns are positions in \C, not regions.}
  \label{fig:registration}
\end{figure}

Registration is thus nothing over and above an instance of \C, and it is always registration
\emph{of} a specified articulation: to say that \R{} registers \F{} without qualification is to leave
the claim incomplete. It is also an actual relation between actual contrasts, not a standing
disposition of \R{} to respond. Figure~\ref{fig:registration} depicts the structure so defined.

The definition attributes to the receiver no understanding, identification, or judgement, since none
of these capacities has been introduced. Where registration is realized physically, it is found in
systems to which nobody would attribute knowledge, and by itself it does not place anything
\enquote{in the logical space of reasons}.\footnote{The phrase is from \textcite[§\,36]{sellars1956};
\textcite[ch.~4]{brandom1994} develops the contrast between such reliable differential
responsiveness and conceptual capacities. In perception, Dretske's non-epistemic seeing, in which an
object is visually differentiated from its surroundings without any belief about it, is a case of
registration in the present sense \autocite{dretske1969}.} Nor does registration move anything from \F{} to \R: the
receiving contrast need neither resemble the articulation nor reproduce its organization, provided
that the dependence obtains. Nor, finally, does it require a representational space, or any
repertoire of alternatives, in \R. Registration relates contrasts that are realized, and nothing in
\R{} need stand for conditions that \R{} does not exhibit.

Registration so defined is supplied by neither relatum in isolation. It comprises three distinct
contributions: from the target, the realized articulation; from the receiver, the contrast between
its conditions; and between them, the dependence that connects the two. Keeping these contributions
distinct prevents registration from being attributed to the target merely because the target is
differentiated, or to the receiver merely because its conditions differ. This does not explain the
emergence of a knowing subject, but it identifies what such an explanation would have to elaborate.

What is registered, on this account, is always a difference within \F, never which way \F{} is. The
core provides for registering that $a$ differs from $b$; it does not provide for registering that
\F{} is $a$ rather than $b$, since that would require $b$ to be an alternative to $a$ rather than a
term beside it. Accounts of information and knowledge built on that richer form presuppose a set of
alternatives among which the actual case is singled out.\footnote{Shannon's theory measures
information as selection from a set of possible messages \autocite{shannon1948}, and Dretske's
informational content is defined by conditional probabilities over such alternatives
\autocite{dretske1981}. In epistemology, relevant-alternatives theories and contrastivism analyse
knowledge as the elimination of alternatives \autocite{dretske1970,lewis1996,schaffer2004};
counterfactual and interventionist analyses of dependence share the same structure
\autocite{lewis1973,woodward2003}.} The core deliberately admits no such set, since it would require
a repertoire of possibilities, held in \F, in \R, or elsewhere, that nothing in the guiding premise
demands. Whether registration of which way \F{} is can be reconstructed from contrasts between actual
terms, without admitting alternatives, is left open.

% -----------------------------------------------------------------------------------
\section{Registration and Order}\label{sec:order}

The core contains no order of any kind. Its contrasts are not sequences, and the conditions of \R{}
are co-present, not successive. \C{} is asymmetric, but the asymmetry of a relation between its
positions is not an ordering of occurrences: that the contrast in \R{} obtains in virtue of the
contrast in \F{} does not place either before the other.

The core therefore cannot describe the receiver as \emph{acquiring} a contrast, \emph{retaining} it,
or \emph{comparing} it with another, and each of these operations requires something
different.\footnote{The three operations closely parallel the synthesis of apprehension,
reproduction, and recognition in the first edition of Kant's Transcendental Deduction
\autocite[A98--A110]{kant1998}, and Husserl's distinction between retention and recollection
\autocite{husserl1991} analyses the second of them. Both treat the operations as presupposing temporal
form; the present account requires only order, of which temporal order is one kind.} Comparison
requires that two receiving contrasts be available together; this is a relation of co-availability,
not of order. Retention requires that one receiver be identified across distinct positions of some
ordering. Acquisition requires that a condition of \R{} in which a contrast is unmatched stand in an
asymmetric relation to one in which it is matched. Only retention and acquisition require order, and
neither follows from (C).

What they require is order in the loosest sense: an asymmetric relation among positions, with no
requirement that any two positions be comparable, and with no metric, no linearity, and no universal
index. Whether this relation must also be transitive, and so form a strict partial order, may be left
open. Temporal succession is one way of realizing such an order; the order of steps in an inference
or a proof is another, which is not temporal. The core is compatible with either and requires
neither.

For the same reason, if receptive dependence is written with a directed symbol, such as an arrow from
\F{} to \R, that symbol should not also serve, without further explanation, to represent order. The
asymmetry of dependence and the asymmetry of an ordering play different roles in the account and
should be kept notationally distinct.

% -----------------------------------------------------------------------------------
\section{Conclusion}\label{sec:concl}

Table~\ref{tab:commitments} summarizes the elements of the account, what each permits, and what each
leaves unestablished. The specified articulation is listed alongside the three commitments, although
it is specified in each application rather than admitted once and for all.

\begin{table}[htbp]
  \centering\small
  \renewcommand{\arraystretch}{1.3}
  \setlength{\tabcolsep}{5pt}
  \captionabove{The elements of the account, with what each permits and what each leaves
                unestablished.}
  \label{tab:commitments}
  \begin{tabularx}{\textwidth}{@{}>{\raggedright\arraybackslash}p{8.8em}
                                    >{\raggedright\arraybackslash}X
                                    >{\raggedright\arraybackslash}X@{}}
    \toprule
    \scshape element & \scshape what it permits & \scshape what it does not establish \\
    \midrule
    (D) Differentiability
      & A differentiation attributable to \F{} itself.
      & Any particular articulation, or its realization. \\
    Specified articulation
      & Terms and the contrast in which they differ.
      & Its own realization, registration, exhaustive partition, or objecthood. \\
    (R) Second relatum
      & Differentiation of \R's own.
      & A knowing subject, or any dependence on \F. \\
    (C) Receptive \mbox{dependence}
      & Registration of a realized articulation; the roles of target and receiver.
      & Which way \F{} is; understanding, representation, order, acquisition, or retention. \\
    \bottomrule
  \end{tabularx}
\end{table}

The reconstruction does not derive everything from differentiability. It shows, rather, why further
commitments are needed, and it distinguishes those commitments from what they make it possible to
define; in particular, neither the second relatum nor its dependence on the target is concealed in
the initial claim about \F. The resulting proposal is that any epistemic relation, in whatever domain
it is realized, requires \emph{differentiation and receptive dependence} between contrasts of actual
terms, with target and receiver distinguished as positions within that dependence. The claim is one
of necessity only: these conditions are required, by the guiding premise, for anything to be
perceived, registered, or claimed about anything, and they are not sufficient for any of these.

The core is minimal, but it is not empty: its commitments exclude several views of what an epistemic
relation presupposes. They exclude, first, the view that all differentiation is supplied by the
receiver, since registration requires a contrast attributable to the field itself
(Section~\ref{sec:diff}). They exclude, second, the view that registration consists in resemblance or
in term-by-term correspondence, since both are symmetric and neither supplies the dependence required
by the second clause of (C) (Section~\ref{sec:dep}). Third, they exclude any conception of
registration as a transfer from target to receiver, since the receiving contrast need neither
resemble nor reproduce the articulation (Section~\ref{sec:reg}). Fourth, they exclude the view that a
repertoire of alternatives, whether held in the field, in the receiver, or elsewhere, is a
precondition of registration, since every contrast in the core is between actual terms
(Sections~\ref{sec:art} and~\ref{sec:reg}). Fifth, they exclude the view that time, or order of any
kind, is such a precondition, since registration has been defined without either
(Section~\ref{sec:order}). Finally, they exclude the view that a knowing subject, concepts, or
representations precede registration, since none of these figures in (D), (R), or (C). None of these
exclusions denies that the structures concerned are present in particular epistemic relations. What
is denied is only that they belong to the core: wherever they are needed, they must enter as further
commitments, stated as such.

Four questions remain open. The first is what grounds the asymmetry of \C{} in a given domain. The
second is whether receptive dependence can itself be explained in terms of more elementary
relations; an arrow whose meaning is left unspecified would not answer it. The third is whether
registration of which way \F{} is can be reconstructed without admitting alternatives. The fourth is
what order acquisition and retention require, and whether anything stronger than an asymmetric
relation is needed. The reconstruction ends at structural registration. Order, memory, recognition,
and representation have not been presupposed in reaching it, and none of them has yet been derived
from it.

\begin{center}
  \vspace{1.2\baselineskip}
  {\large ⁂}
\end{center}

\printbibliography[heading=bibintoc,title={References}]

\end{document}
